% mode: ini-etex
% What §422 and §360 show of an argument read by a capture loop: everything before the token that stopped it, from each level it came
% from, and nothing after.  align_state has to come back right too, so each runs in a cell, where a wrong one would stop & from firing.
\input prelude
\tracingmacros=1 \errorcontextlines=100
\def\a#1{[#1]}
\def\c#1.{[#1]}
\def\h#1abc{[#1]}
\def\f#1#2.{[#1|#2]}
\def\D#1{\noindent\par\message{^^J#1:}}
% Run in a cell, at align_state 0, and at the top level.
\def\cell#1{\setbox0\vbox{\halign{##&##\cr x#1&y\cr}}\showbox0 \setbox0\hbox{#1}\showbox0 \noindent\par}
\long\def\run#1{\def\q{#1}\expandafter\cell\expandafter{\q}\noindent\par}

\D{par ends a non-long argument}
\run{\a\par z}
\run{\c ab\par z.}
\run{\c{ab}cd\par z.}
\run{\c ab{cd\par e}f.}
\run{\c ab{c{d{e\par f}g}h}i.}
\run{\a{ab\par c}}
\run{\a{a{b{c\par d}e}f}}
\run{\f{a}bcd\par z.}
\run{\h aba\par c abc}
\run{\h ab\par abc}
\run{\h a\par b abc}

\D{from several levels}
\def\ta{\c a}
\long\def\tb{b\par z.}
\long\def\tc{\ta\tb}
\cell{\ta b\par z.}
\cell{\tc}
\def\td#1#{\a}
\cell{\td{ab\par c}}
\def\te{\td{ab\tf}}
\long\def\tf{cd\par e}
\cell{\te}

\D{in an argument read in place}
\def\ra#1{\c x#1y.}
\long\def\rl#1{\ra{#1}}
\long\def\rz#1{\expandafter\cell\expandafter{\rl{#1}}}
\cell{\ra{p\par q}}
\long\def\rb#1{\a{x#1y}}
\cell{\rb{p\par q}}
\cell{\rb{\par}}
\cell{\rb{{p\par q}}}
\long\def\rc#1{\a{x#1y\par z}}
\cell{\rc{p}}
\long\def\rd#1#2{\a{x#1y#2z}}
\cell{\rd{p}{q\par r}}
\cell{\rd{p\par}{q}}
\long\def\re#1{\c x#1y.}
\cell{\re{p\par q}}
\cell{\re{\par}}
\long\def\rf#1{\h x#1y abc}
\cell{\rf{p\par q}}
\cell{\rf{a\par}}
\cell{\rf{ab\par}}

\D{a long macro lets it through, a macro it calls does not}
\long\def\la#1{\a{#1}}
\cell{\la{p\par q}}
\long\def\lb#1{\a#1}
\cell{\lb{\par}}
\long\def\lc#1{\c#1.}
\cell{\lc{p\par q}}
\long\def\ld#1{[#1]}
\long\def\le#1.{[#1]}
\long\def\lf#1{\ld{#1}\le#1.}
\cell{\lf{p\par q}}

\D{the state after an error: an extra right brace makes the argument's long_state CALL}
\long\def\ed#1.{[#1]}
\def\ee{\ed a}
\ee}b\par c.
\noindent\par
\ee}b.
\noindent\par
\def\ef#1{\a#1}
\ef}
\noindent\par
\def\eg{\c ab}
\eg}z.
\noindent\par
\def\eh{\f{a}}
\eh}z.
\noindent\par
\def\ei{\h ab}
\ei}abc
\noindent\par
\def\ej{\a  }
\ej}z
\noindent\par
% The right brace in the same list: a group opened in front keeps the text balanced.  Not in a cell: §421 leaves align_state 1 too high, on any path.
\def\inh#1{\setbox0\hbox{#1}\showbox0 \noindent\par}
\def\q{{\ed a}z}\inh{\q}
\def\q{{\c a}z}\inh{\q}
\def\q{{\a}z}\inh{\q}
\def\q{{\a }z}\inh{\q}
\def\q{{\a{a}}z}\inh{\q}
\def\q{{\c{a}}z}\inh{\q}
\def\q{{\h ab}z}\inh{\q}
\def\q{{\f{a}}z}\inh{\q}
\def\q#1{{\c a#1}z}\inh{\q{}}\inh{\q{b}}
\def\q#1{{\a#1}z}\inh{\q{}}
\def\q#1{{\c#1}z}\inh{\q{{a}}}

\D{no match of a prefix}
\def\pa abc#1{[#1]}
\def\q{\pa abd}\q
\def\q{\pa ab\par}\q
\def\q{\pa ab{}}\q
\noindent\par

\D{a runaway with an outer macro}
\def\O{plain}
\def\oa{\c ab\O cd.}
\def\ob{\a{ab\O cd}}
\def\oc#1{\a{ab#1cd}}
\def\od#1{\c ab#1cd.}
\def\oe#1{\oP\a{ab#1cd}}
\let\oP=\relax
\oc{\O}\od{\O}\oe{\O}
\def\oP{\outer\def\O{outer}}
\oe{\O}
\noindent\par
\oa
\noindent\par
\ob
\noindent\par
\message{^^Jat the end of the document, no delimiter before the end of the file}
\def\ga#1.{[#1]}
\def\gb{\ga ab}
\gb
